%=============================================================
%  Module I.4 / L09 — The Hydrogen Atom
%  Series I > Module I.4 > Lecture L09
%  Duration: 90 minutes  |  All tiers (T1–T5)
%=============================================================
\renewcommand{\lectureCode}{I.4 / L09}

\section*{L09 — The Hydrogen Atom}
\label{sec:L09}
\addcontentsline{toc}{section}{L09 — The Hydrogen Atom}

%--- Lecture header
\begin{tcolorbox}[lecturebox, title={Module I.4 / L09 — Metadata}]
\begin{tabular}{@{}ll@{}}
\textbf{Series:}    & I — Introductory Quantum Mechanics \\
\textbf{Module:}    & I.4 — Paradigm Physical Systems \\
\textbf{Lecture:}   & L09 — The Hydrogen Atom \\
\textbf{Duration:}  & 90 minutes \\
\textbf{Tiers:}     & T1 (HS) · T2 (BegUG) · T3 (AdvUG) · T4 (MSc) · T5 (PhD) \\
\textbf{Preceding:} & I.4 / L08 — Spin–Orbit Coupling and Fine Structure \\
\textbf{Following:} & I.4 / L10 — Scattering from a Potential Barrier \\
\textbf{Prerequisites:} & Modules I.1–I.3 complete \\
\end{tabular}

\medskip
\textbf{Key concepts:} Coulomb potential, radial equation, Laguerre polynomials, SO(4) symmetry, Runge–Lenz vector
\end{tcolorbox}

%------------------------------------------------------------
\subsection*{L09.0 — Learning Objectives (per tier)}
\label{sec:L09.0}

\tiertag{tier1col}{Tier 1 — HS}\quad
Identify the main physical phenomenon studied in this lecture; describe it in
non-mathematical terms; state one real-world application.

\tiertag{tier2col}{Tier 2 — BegUG}\quad
Apply the key equations to compute numerical results; verify consistency with
known limits; translate between physical and mathematical descriptions.

\tiertag{tier3col}{Tier 3 — AdvUG}\quad
Derive central results from first principles; prove key identities; identify
the mathematical structure underlying the physical phenomenon.

\tiertag{tier4col}{Tier 4 — MSc}\quad
Analyse formal extensions; connect to symmetry principles; generalise to
related systems; identify the role of this lecture in the broader programme.

\tiertag{tier5col}{Tier 5 — PhD}\quad
Establish rigorous mathematical foundations; identify open research problems;
connect to modern literature; critique standard textbook treatments.

%--- Lecture content -----------------------------------------

\subsection*{L09.1 — Separation in Spherical Coordinates}
$\hat{H} = -\frac{\hbar^2}{2m}\nabla^2 - \frac{e^2}{4\pi\epsilon_0 r}$;
$\psi_{nlm}(r,\theta,\phi)=R_{nl}(r)Y_l^m(\theta,\phi)$.
Radial substitution $u_{nl}=rR_{nl}$; effective potential
$V_{\rm eff}(r) = -\frac{e^2}{4\pi\epsilon_0 r} + \frac{\hbar^2 l(l+1)}{2mr^2}$.

\subsection*{L09.2 — Energy Spectrum}
\begin{tcolorbox}[lecturebox, title={Hydrogen Spectrum}]
\begin{equation}
  E_n = -\frac{m_e e^4}{2(4\pi\epsilon_0)^2\hbar^2}\frac{1}{n^2}
       = -\frac{13.6\,\text{eV}}{n^2}, \quad n=1,2,3,\ldots
  \label{eq:L09.1}
\end{equation}
Quantum numbers: $n\geq1$; $l=0,\ldots,n-1$; $m=-l,\ldots,l$.
Degeneracy: $n^2$ (ignoring spin), $2n^2$ with spin.
\end{tcolorbox}

\subsection*{L09.3 — Hidden SO(4) Symmetry}
\begin{tcolorbox}[lecturebox, title={Runge–Lenz Vector}]
$\hat{\mathbf{A}} = \frac{1}{2m}(\hat{\mathbf{p}}\times\hat{\mathbf{L}}-\hat{\mathbf{L}}\times\hat{\mathbf{p}})
  - \frac{e^2}{4\pi\epsilon_0}\frac{\hat{\mathbf{r}}}{r}$;
$[\hat{H},\hat{\mathbf{A}}]=0$. Together with $\hat{\mathbf{L}}$, forms $\mathfrak{so}(4)$ algebra
$\Rightarrow$ $1/n^2$ spectrum from algebra alone.
\end{tcolorbox}

\subsection*{L09.4 — Expectation Values}
$\langle r\rangle_{nl} = \frac{a_0}{2Z}[3n^2-l(l+1)]$;
$\langle r^{-1}\rangle = Z/n^2a_0$;
$\langle r^{-3}\rangle_{nl} = Z^3/[n^3a_0^3 l(l+\tfrac{1}{2})(l+1)]$ (Kramers–Pasternack).


%------------------------------------------------------------
\subsection*{L09.C — Connections}
\label{sec:L09.C}
\textbf{Backward:} I.4/L08 — Spin–Orbit Coupling and Fine Structure and Modules I.1–I.3.
\textbf{Forward:} I.4/L10 — Scattering from a Potential Barrier builds directly on the formalism developed here.
\textbf{Cross-module:} Series~II modules extend these results to many-body and
advanced settings; Series~III applies them to molecular electronic structure.

%=============================================================
%  ASSESSMENT BUNDLE — L09
%=============================================================
\newpage
\section*{L09 — Assessment Artifact Bundle}

\begin{tcolorbox}[ugconceptbox, title={SCQU · L09 · Tiers 1–3 · 15–20 min}]
\textit{Ten simple concept questions (mix MC/TF/short) targeting the core results
of The Hydrogen Atom. See artifact file \texttt{L09_artifacts.tex} for full questions.}

\textbf{Rubric:} 2 pts each (1.5 correct + 0.5 justification). Total: 20 pts.
\end{tcolorbox}

\begin{tcolorbox}[ugconceptbox, title={CCQU · L09 · Tiers 2–3 · 25–35 min}]
\textit{Ten complex concept questions: ``explain why,'' ``what would change if,''
``identify the flaw.'' Full questions in \texttt{L09_artifacts.tex}.}

\textbf{Rubric:} 2 pts each (1 key point + 1 argument). Total: 20 pts.
\end{tcolorbox}

\begin{tcolorbox}[ugprobbox, title={SPSU · L09 · Tiers 1–3 · 60–90 min}]
\textit{Ten problems (Part A: mechanics, Part B: interpretation, Part C: translation).
Full problems in \texttt{L09_artifacts.tex}.}

\textbf{Rubric:} Result 60\%, Method 30\%, Interpretation 10\%.
\end{tcolorbox}

\begin{tcolorbox}[ugprobbox, title={CPSU · L09 · Tier 3 · 3–6 hrs (4-layer)}]
\textit{Five multi-part derivation problems with four-layer scaffolding.
Layers 1–4 in \texttt{L09_ex01\_problem.tex}, \texttt{hints}, \texttt{adv\_hints}, \texttt{solution}.}

\textbf{Rubric:} Assumptions 15\%, Proof structure 45\%, Algebra 25\%, Insight 15\%.
\end{tcolorbox}

\begin{tcolorbox}[ugprobbox, title={SPJU + CPJU · L09 · Tiers 2–3}]
\textit{Simple project (2–6 hrs): structured computational/analytical task.
Complex project (8–15 hrs): open-ended synthesis.
Full specs in \texttt{L09_artifacts.tex}.}
\end{tcolorbox}

\begin{tcolorbox}[gradconceptbox, title={SCQG + CCQG · L09 · Tiers 4–5}]
\textit{Ten simple + ten complex concept questions at graduate level (rigour, domain
conditions, named theorems). Full questions in \texttt{L09_artifacts.tex}.}
\end{tcolorbox}

\begin{tcolorbox}[gradprobbox, title={SPSG + CPSG · L09 · Tiers 4–5}]
\textit{Ten simple + five complex graduate problems.
CPSG with four-layer format. See \texttt{L09_artifacts.tex}.}
\end{tcolorbox}

\begin{tcolorbox}[gradprobbox, title={SPJG (×5) + CPJG (×5) · L09 · Tiers 4–5}]
\textit{Five simple (3–8 hrs each) and five complex (10–20 hrs each) graduate projects.
Full specs in \texttt{L09_artifacts.tex}.}
\end{tcolorbox}

\begin{tcolorbox}[researchbox, title={WRQ (×5) + ORQ (×5) · L09 · Tiers 4–5}]
\textit{Five well-defined research questions (5–10 hrs each) and five open-ended
research questions (10–20 hrs each). Full prompts in \texttt{L09_artifacts.tex}.}
\end{tcolorbox}
